\chapter{Dynamics and Decision Windows}
\label{ch:dynamics}

\section{Volitional Energy}

\begin{definition}[Volitional energy]
Volitional energy $\Will(t)$ is the limited capacity available at time $t$ to
interrupt the current behavioural frame and initiate stabilizing action toward
a target frame.
\end{definition}

\begin{definition}[Leverage]
The effective leverage of free will at time $t$ is
\[
  \Leverage(t) = \frac{\Will(t)}{\DV(t)}.
\]
\end{definition}

\begin{theorem}[Decision-window principle]
Free will has maximum strategic effect not when effort is most intense, but
when the activation barrier is locally minimal relative to available
volitional energy.
\end{theorem}

\begin{proof}[Proof sketch]
If leverage is represented by $\Will(t)/\DV(t)$, then leverage rises when
$\Will(t)$ increases or $\DV(t)$ decreases. In ordinary life, large increases
in available will are unreliable, but temporary drops in barrier are common:
after waking, after finishing a task, upon entering a new environment, or
immediately after a sharp emotional update. Therefore the most reliable
strategy is to act during barrier troughs rather than waiting for heroic
motivation.
\end{proof}

\section{Decision Windows}

\begin{definition}[Decision window]
A decision window is a short interval during which the current state has lost
part of its stabilizing force while the next state has not yet been rejected.
\end{definition}

Common windows include:
\begin{itemize}[nosep]
  \item waking before distraction is loaded;
  \item the minute after finishing a scheduled block;
  \item physical arrival at a task-specific location;
  \item social commitments that force a new frame to begin.
\end{itemize}

\section{Failure Modes}

\begin{lemma}[Late-response failure]
If the agent waits until inertia has fully rebuilt after a window closes, the
same target action requires strictly greater effort than during the window.
\end{lemma}

\begin{discussion}
This lemma explains a common misunderstanding. People often believe they have
``decided'' to act because they endorsed a future action verbally. But if they
do not move inside the decision window, the state soon becomes expensive again
and the earlier verbal decision has almost no mechanical value.
\end{discussion}

\section{A Comparative Map of Windows}

\begin{center}
\begin{tabularx}{\textwidth}{@{}l l X@{}}
\toprule
Window type & Typical duration & Strategic use \\
\midrule
Morning reset & 5--20 minutes & Inject a pre-written plan before default inputs arrive \\
Task boundary & 1--5 minutes & Move directly into the next prepared task while switching cost is already being paid \\
Spatial transition & Variable & Attach target behaviour to places that make old cues less available \\
Social trigger & Variable & Convert personal intention into a scheduled external obligation \\
\bottomrule
\end{tabularx}
\end{center}

\section{Recommended Reading}

\begin{readingbox}
\textbf{Foundation}
\begin{itemize}[nosep]
  \item Steven Strogatz, \emph{Nonlinear Dynamics and Chaos}.
  \item Eugene Izhikevich, \emph{Dynamical Systems in Neuroscience}.
\end{itemize}
\textbf{Modern Research}
\begin{itemize}[nosep]
  \item Nature Neuroscience work on predictable bifurcation dynamics during
        falling asleep.
  \item Review work on metastability in brain dynamics.
\end{itemize}
\textbf{Frontier}
\begin{itemize}[nosep]
  \item Use frontier network models only after the basic decision-window logic
        is clear.
\end{itemize}
\end{readingbox}

\begin{summarybox}
The second chapter explains why raw effort is a poor operating principle.
Timing dominates intensity. The business value of the model comes from locating
repeatable low-barrier windows and treating them as scarce assets.
\end{summarybox}
